\documentclass[12pt,a4paper]{article}
\usepackage{iftex}
\ifPDFTeX
  \usepackage[T2A]{fontenc}
  \usepackage[utf8]{inputenc}
\else
  \usepackage{fontspec}
  \IfFontExistsTF{Times New Roman}{\setmainfont{Times New Roman}}{\setmainfont{DejaVu Serif}}
  \IfFontExistsTF{Arial}{\setsansfont{Arial}}{\setsansfont{DejaVu Sans}}
  \IfFontExistsTF{Consolas}{\setmonofont{Consolas}}{\setmonofont{DejaVu Sans Mono}}
\fi
\usepackage[ukrainian]{babel}
\usepackage{microtype,geometry,setspace,enumitem,titlesec,booktabs,array,tabularx,longtable}
\usepackage{xcolor,listings,graphicx,float,fancyhdr,lastpage,hyperref,tikz}
\usetikzlibrary{arrows.meta,positioning}
\geometry{left=25mm,right=15mm,top=20mm,bottom=20mm}
\setstretch{1.15}
\setlength{\parindent}{1.25cm}
\setlength{\parskip}{0.18em}
\setlist[itemize]{leftmargin=1.8em,itemsep=0.2em,topsep=0.25em}
\setlist[enumerate]{leftmargin=1.9em,itemsep=0.2em,topsep=0.25em}
\setlength{\emergencystretch}{4em}
\renewcommand{\arraystretch}{1.18}
\definecolor{primary}{HTML}{163629}
\definecolor{accent}{HTML}{287857}
\definecolor{soft}{HTML}{F1F6F0}
\definecolor{codebg}{HTML}{F6F8FA}
\definecolor{codegreen}{HTML}{197A45}
\definecolor{codegray}{HTML}{5F6B7A}
\hypersetup{colorlinks=true,linkcolor=primary,urlcolor=accent,citecolor=primary,
  pdftitle={Звіт до лабораторної роботи з Web-технологій},pdfauthor={Кашинський Владислав}}
\titleformat{\section}{\large\bfseries\color{primary}}{\makebox[2.4em][l]{\thesection.}}{0pt}{}
\titleformat{\subsection}{\normalsize\bfseries\color{primary}}{\makebox[3.0em][l]{\thesubsection.}}{0pt}{}
\titleformat{\subsubsection}{\normalsize\bfseries\color{accent}}{\makebox[4.1em][l]{\thesubsubsection.}}{0pt}{}
\pagestyle{fancy}
\fancyhf{}
\fancyhead[L]{\small Web-технології}
\fancyhead[R]{\small StudySpace}
\fancyfoot[C]{\small \thepage\ з \pageref{LastPage}}
\renewcommand{\headrulewidth}{0.4pt}
\setlength{\headheight}{15pt}
\lstdefinestyle{studyspace}{backgroundcolor=\color{codebg},basicstyle=\ttfamily\small,
  keywordstyle=\color{accent}\bfseries,commentstyle=\color{codegreen},stringstyle=\color{primary},
  numbers=left,numberstyle=\tiny\color{codegray},numbersep=8pt,frame=single,
  rulecolor=\color{codegray!30},breaklines=true,showstringspaces=false,tabsize=2}
\lstset{style=studyspace}
\newcommand{\code}[1]{\path{#1}}
\newcommand{\reportfigure}[3]{\begin{figure}[H]\centering
  \IfFileExists{#1}{\includegraphics[width=#2\textwidth,height=0.70\textheight,keepaspectratio]{#1}}
  {\fbox{\parbox[c][5cm][c]{0.88\textwidth}{\centering Знімок екрана відсутній.}}}
  \caption{#3}\end{figure}}
\newcommand{\taskcard}[5]{\subsubsection{#1: #2 (#3 б.)}
  \textbf{Мета.} #4

  \textbf{Реалізація і перевірка.} #5}

\begin{document}
\thispagestyle{empty}
\begin{center}
{\large Київський національний університет імені Тараса Шевченка}\\[0.5em]
{\large Факультет інформаційних технологій}\\
{\large Кафедра програмних систем і технологій}\\[4.6cm]
{\Large\bfseries ЗВІТ}\\[0.5em]
{\large до лабораторної роботи}\\[0.25em]
{\large з дисципліни «Web-технології»}\\[0.9em]
{\Large\bfseries Розроблення платформи спільної роботи StudySpace}
\end{center}
\vfill
\begin{flushright}\begin{minipage}{0.48\textwidth}
Виконав:\\ студент групи ТТП-32\\ \textbf{Кашинський Владислав}\\[1.2em]
\end{minipage}\end{flushright}
\vfill
\begin{center}Київ -- 2026\end{center}

\newpage
\section*{Анотація}
\addcontentsline{toc}{section}{Анотація}
У лабораторній роботі спроєктовано та реалізовано StudySpace, навчальну
web-платформу для організації робочих просторів, документів, файлів і
спільної роботи. Пов'язані вимоги об'єднано в один Git-монорепозиторій,
а технологічно незалежні частини оформлено як окремі застосунки. Рішення
містить ASP.NET Core MVC/API, React SPA, Azure Functions, RabbitMQ Worker,
MCP-сервер, SQLite, SignalR, SSE, WebSocket, Lucene.NET,
Leaflet/OpenStreetMap, Chart.js та автоматизовані тести.

Обрано 36 пунктів завдання сумарною вартістю 104 бали. Під час перевірки
підтверджено форматування і збірку .NET-рішення, 20 модульних та
інтеграційних тестів, 5 browser E2E-сценаріїв, production-збірку React,
REST CRUD, пагінацію, кеш, пошук, SSE і MCP. Зовнішні cloud-інтеграції
залишено вимкненими до додавання особистих облікових даних студента.

\textbf{Ключові слова:} ASP.NET Core, MVC, REST API, React, SignalR,
WebSocket, SSE, SQLite, Entity Framework Core, RabbitMQ, Azure Functions,
MCP, Playwright, Docker.

\newpage
\tableofcontents
\newpage

\section{Постановка задачі}
\subsection{Мета роботи}
Мета роботи полягає у проходженні повного циклу створення web-системи: від моделі даних і
CRUD до REST API, браузерного інтерфейсу, авторизації, realtime-обміну,
інтеграцій, тестування та підготовки розгортання. Практичний результат має
працювати локально й бути готовим до подальшої публікації.

\subsection{Предметна область}
StudySpace зберігає навчальні матеріали у робочих просторах. Простір має
назву, опис, тариф, координати та учасників. У ньому створюються документи зі
статусом, summary, обкладинкою, впорядкованими блоками та файлами. Учасники
мають ролі Owner, Editor або Viewer. Premium відкриває дошку, помічника та
розширені можливості.

\subsection{Вибір завдань і розрахунок}
Сума балів за реалізовані завдання становить 104. За формулою з умови
номінальний результат:
\[60\cdot\frac{104}{100}=62{,}4.\]

\begin{longtable}{>{\bfseries}p{0.25\textwidth} p{0.57\textwidth} r}
\toprule Етапи & Практичний результат & Бали \\
\midrule\endfirsthead
\toprule Етапи & Практичний результат & Бали \\
\midrule\endhead
C1, C5, QA1, B1 & MVC CRUD, файли, REST, Swagger/Postman, пагінація & 10 \\
C2, C3, C4 & Azure-ready deployment, GitHub Actions, Docker/Compose & 5 \\
C6, C7, F1, F2, QA2 & Landing, SEO, Chart.js, AJAX, E2E & 13 \\
C8, B2, C9, B4 & Карта, webhook/Telegram, telemetry, Lucene & 9 \\
C10, C11 & Polling, long polling, WebSocket, SignalR & 8 \\
C12, F3, C13, C14 & Block editor, drag-and-drop, realtime, presence & 9 \\
C15 & File API, DataTransfer, FileReader & 3 \\
C16, B5, B6, QA3 & OAuth, ролі, multi-tenancy, plans, auth tests & 12 \\
F5 & Окремий React + TypeScript SPA & 3 \\
B7 & Окремий Azure Functions API & 3 \\
C18, F6 & Shared board, Pointer Events, Canvas 2D & 9 \\
B8, B9, C19 & RabbitMQ, memory cache, Azure Blob/local storage & 10 \\
C21, C22 & SSE-помічник і MCP tools & 10 \\
\midrule \multicolumn{2}{r}{\textbf{Разом}} & \textbf{104} \\
\bottomrule
\end{longtable}

\section{Поділ рішення на проєкти}
\subsection{Чому монорепозиторій}
Усі частини працюють з однією предметною областю, тому їх збережено в одному
репозиторії. Це дає спільну історію, одну команду збірки й відсутність
дублювання моделей. Водночас кожен незалежний процес є окремим проєктом.
Результат: \textbf{один репозиторій}, \textbf{5 запускних застосунків} і
\textbf{10 технічних проєктів}.

\subsection{Запускні застосунки}
\begin{table}[H]\centering\small
\begin{tabularx}{\textwidth}{>{\bfseries}p{0.25\textwidth} p{0.28\textwidth} X}
\toprule Застосунок & Розташування & Роль \\
\midrule
StudySpace Web & \code{src/StudySpace.Web} & MVC, REST, SignalR, SSE, WebSocket, UI \\
StudySpace SPA & \code{src/studyspace-spa} & React/TypeScript REST-клієнт \\
StudySpace Functions & \code{src/StudySpace.Functions} & Serverless health і digest \\
StudySpace Worker & \code{src/StudySpace.Worker} & RabbitMQ consumer \\
StudySpace MCP & \code{src/StudySpace.Mcp} & Stdio JSON-RPC/MCP tools \\
\bottomrule
\end{tabularx}\caption{П'ять окремо запущуваних частин}\end{table}

\subsection{Усі технічні проєкти}
\begin{longtable}{>{\bfseries}p{0.25\textwidth} p{0.50\textwidth} r}
\toprule Проєкт & Призначення & Тип \\
\midrule\endfirsthead
\toprule Проєкт & Призначення & Тип \\
\midrule\endhead
StudySpace.Core & Сутності й enum предметної області & .NET library \\
StudySpace.Data & DbContext, EF Core, SQLite, migration, seed & .NET library \\
StudySpace.Web & MVC, API, realtime, UI, integrations & .NET web \\
StudySpace.Functions & Ізольовані Azure Functions & .NET functions \\
StudySpace.Worker & RabbitMQ background service & .NET worker \\
StudySpace.Mcp & MCP stdio server & .NET console \\
StudySpace.\newline UnitTests & Доменні unit-тести & .NET tests \\
StudySpace.\newline IntegrationTests & API й authorization tests & .NET tests \\
studyspace-spa & React, TypeScript, Vite & Node.js \\
e2e & Playwright Chromium scenarios & Node.js \\
\bottomrule\end{longtable}

\subsection{Структура каталогів}
\begin{lstlisting}[language={},caption={Скорочена структура}]
StudySpace/
|-- src/
|   |-- StudySpace.Core/        |-- StudySpace.Data/
|   |-- StudySpace.Web/         |-- StudySpace.Functions/
|   |-- StudySpace.Worker/      |-- StudySpace.Mcp/
|   `-- studyspace-spa/
|-- tests/
|   |-- StudySpace.UnitTests/   |-- StudySpace.IntegrationTests/
|   `-- e2e/
|-- docs/stages/                |-- report/WebTechnologiesLab/
|-- .github/workflows/ci.yml    |-- compose.yaml
`-- StudySpace.slnx
\end{lstlisting}

\section{Архітектура і модель даних}
\subsection{Архітектурна схема}
\begin{figure}[H]\centering
\begin{tikzpicture}[node distance=9mm and 10mm,
box/.style={draw=primary,rounded corners,fill=soft,align=center,minimum height=10mm,text width=30mm},
line/.style={-{Latex[length=2mm]},thick,color=primary}]
\node[box] (browser) {MVC UI\\Browser};
\node[box,above=of browser] (spa) {React SPA};
\node[box,right=of browser] (web) {StudySpace.Web\\MVC + REST + realtime};
\node[box,below=of web] (data) {Core + Data\\EF Core + SQLite};
\node[box,right=of web] (queue) {RabbitMQ\\Worker};
\node[box,below=of queue] (storage) {Azure Blob\\або local files};
\node[box,above=of queue] (mcp) {MCP Server};
\node[box,right=of queue] (func) {Azure\\Functions};
\draw[line] (browser)--(web); \draw[line] (spa)--(web); \draw[line] (web)--(data);
\draw[line] (web)--(queue); \draw[line] (web)--(storage); \draw[line] (mcp)--(web);
\draw[line] (func)--(web);
\end{tikzpicture}\caption{Взаємодія компонентів StudySpace}\end{figure}

Основна точка входу налаштовує dependency injection, EF Core, cookies,
опційний Google OAuth, Swagger, SignalR, Application Insights і middleware.
Контролери перед читанням або зміною даних виконують tenant-aware перевірку.

\subsection{Сутності}
\begin{table}[H]\centering\small
\begin{tabularx}{\textwidth}{>{\bfseries}p{0.25\textwidth} X}
\toprule Сутність & Основні поля та зв'язки \\
\midrule
Workspace & Id, Name, Description, Plan, coordinates; Documents і Members \\
Document & WorkspaceId, Title, Summary, CoverImageUrl, Status; Blocks, Attachments \\
DocumentBlock & DocumentId, Type, Content, Position \\
Attachment & DocumentId, FileName, ContentType, Size, Url \\
WorkspaceMember & WorkspaceId, UserId, DisplayName, Role \\
\bottomrule\end{tabularx}\caption{Модель предметної області}\end{table}

SQLite обрано для запуску без окремого сервера. EF Core створює схему через
міграції. Seed ідемпотентний: демонстраційні дані додаються лише у порожню БД.

\subsection{Ключові маршрути}
\begin{table}[H]\centering\small
\begin{tabularx}{\textwidth}{>{\ttfamily}p{0.34\textwidth} X}
\toprule Маршрут & Призначення \\
\midrule
/ & dashboard \\
/Workspaces & список і CRUD просторів \\
/Documents/Details/\{id\} & документ, блоки, файли \\
/Documents/Editor/\{id\} & realtime editor \\
/Files & бібліотека вкладень \\
/Insights & статистика й Chart.js \\
/Connectivity & три способи оновлення \\
/Board & Canvas-дошка \\
/Map & Leaflet/OpenStreetMap \\
/swagger & REST документація \\
/api/assistant/stream & SSE-помічник \\
\bottomrule\end{tabularx}\caption{Ключові маршрути}\end{table}

\section{Послідовність роботи системи}
Знімки зроблено під час фактичного локального запуску
\url{http://localhost:5036}; це не макети. Інтерфейс застосунку локалізовано
українською мовою.

\subsection{Крок 1. Запуск і dashboard}
Під час запуску застосовуються міграції та seed. Dashboard показує кількість
просторів, документів, опублікованих матеріалів і файлів.
\reportfigure{images/01-home.png}{0.96}{Головна сторінка після локального запуску}

\subsection{Крок 2. Робочі простори}
Список фільтрується за поточним користувачем. Картка містить тариф, опис і
кількість документів; кнопка створення простору відкриває валідовану MVC-форму.
\reportfigure{images/02-workspaces.png}{0.96}{Список доступних робочих просторів}

\subsection{Крок 3. Конкретний workspace}
Власник бачить документи й дозволені дії: створення, налаштування, спільну
дошку та асистента.
\reportfigure{images/03-workspace-details.png}{0.96}{Деталі робочого простору}

\subsection{Крок 4. Створення документа}
Форма приймає назву, короткий опис, URL обкладинки й статус. Сервер перевіряє
модель, доступ та унікальність, після POST виконує redirect і, за наявності
налаштувань, telemetry, RabbitMQ та webhook.
\reportfigure{images/04-create-document.png}{0.96}{MVC-форма створення документа}

\subsection{Крок 5. Деталі документа}
Сторінка об'єднує відомості, блоки та вкладення. Звідси відкриваються редактор,
редагування та завантаження файлів.
\reportfigure{images/05-document-details.png}{0.96}{Документ із контентом і формою вкладень}

\subsection{Крок 6. Блоковий редактор}
Заголовок, абзац, список, код і цитата зберігаються через REST після короткої
затримки. Перетягування змінює порядок, SignalR передає зміни та фокус.
\reportfigure{images/06-block-editor.png}{0.96}{Блоковий realtime-редактор}

\subsection{Крок 7. Файли}
Зона завантаження читає DataTransfer.files, FileReader будує локальний перегляд, сервер
перевіряє розширення й тарифний ліміт.
\reportfigure{images/07-files.png}{0.96}{Бібліотека до завантаження}

\subsection{Крок 8. Аналітика}
Сторінка аналітики рахує дані з БД; Chart.js будує діаграму статусів.
\reportfigure{images/09-insights.png}{0.96}{Аналітика з Chart.js}

\subsection{Крок 9. REST API}
Swagger дозволяє виконати CRUD, пагінацію, кешування та пошук. Інтеграційні тести
перевіряють HTTP-коди та JSON-контракт.
\reportfigure{images/10-swagger.png}{0.96}{Swagger/OpenAPI endpoints}

\subsection{Крок 10. Мережеві механізми}
Один сигнал паралельно читається polling, long polling і WebSocket. Екран
рахує запити, повідомлення й затримку.
\reportfigure{images/11-connectivity.png}{0.96}{Порівняння трьох transport}

\subsection{Крок 11. Спільна дошка}
Canvas 2D реагує на Pointer Events, координати нормалізуються, штрихи та очищення
синхронізуються SignalR.
\reportfigure{images/12-board.png}{0.96}{Спільна дошка в режимі realtime}

\subsection{Крок 12. Карта}
Leaflet завантажує маркери доступних просторів, OpenStreetMap надає карту, а фільтр
за тарифом повторно запитує API.

\subsection{Крок 13. SSE-помічник}
Помічник передає етапи пошуку й відповідь частинами через SSE. Без API key працює
локальний контекстний варіант.
\reportfigure{images/14-assistant.png}{0.96}{Асистент робочого простору}

\section{Детальний опис 36 завдань}
\subsection{Common}
\taskcard{C1}{MVC і модель даних}{3}{Створити пов'язаний CRUD із файлом.}
{У Core створено п'ять сутностей, у Data створено DbContext, міграції й seed.
WorkspacesController і DocumentsController реалізують CRUD, AttachmentStorage
файли, CoverImageUrl містить зовнішнє посилання. Перевірено всі операції та
повторний запуск без дублювання seed.}

\taskcard{C2}{Azure deployment}{1}{Підготувати cloud-публікацію.}
{Workflow виконує Release publish і conditional azure/webapps-deploy. Назва
ресурсу та publish profile винесені у vars/secrets. Release build пройшов;
живий URL потребує особистої Azure subscription.}

\taskcard{C3}{Continuous integration}{2}{Автоматизувати перевірки змін.}
{\code{ci.yml} має jobs для restore, format, build, tests, container,
frontend та E2E. TRX і Playwright report публікуються як artifacts. Ключові
команди відтворено локально.}

\taskcard{C4}{Контейнеризація}{2}{Описати web, worker і broker.}
{Multi-stage Dockerfiles створюють Web і Worker images, Compose під'єднує
RabbitMQ і volumes. Файли перевірено статично та збіркою проєктів; Docker
Desktop відсутній, тому runtime Compose слід перевірити після інсталяції.}

\taskcard{C5}{REST API}{3}{Надати CRUD документів.}
{API повертає DTO. POST дає 201+Location, PUT 200, DELETE 204; є 400, 403,
404 і 409. Lifecycle пройшов integration і ручний smoke test, тимчасовий
запис видалено.}

\taskcard{C6}{Адаптивний UI}{3}{Створити зрозумілий dashboard.}
{Bootstrap, responsive grid, cards і forms формують єдину design system.
Навігація об'єднує модулі. Landing і переходи перевірено Playwright.}

\taskcard{C7}{SEO}{2}{Додати засоби індексації.}
{Є description, canonical, Open Graph, dynamic robots.txt і sitemap.xml.
Routes покриті E2E; зовнішній SEO audit можливий після deployment.}

\taskcard{C8}{Карта}{3}{Відобразити геодані й фільтр.}
{Map/API повертають доступні координати; Leaflet показує OpenStreetMap і
markers, plan selector фільтрує дані. Marker перевірено E2E і скріншотом.}

\taskcard{C9}{Application Insights}{2}{Підготувати telemetry.}
{SDK активується лише з connection string; document create/update фіксують
custom events. Без секрету використовується безпечний no-op; Portal потребує
власного ресурсу.}

\taskcard{C10}{Polling, long polling, WebSocket}{4}{Практично порівняти transport.}
{ConnectivityController реалізує state, очікування нової версії та WebSocket.
Клієнт рахує запити, повідомлення й latency. Локально всі три канали працювали.}

\taskcard{C11}{SignalR}{4}{Додати високорівневу realtime-взаємодію.}
{CollaborationHub має groups за document/board id і передає blocks, focus,
strokes та clear. Editor і Board показали Realtime connected.}

\taskcard{C12}{Блоковий редактор}{3}{Створити Notion-подібний контент.}
{DocumentBlock має type/content/position. API та UI підтримують п'ять типів,
debounced save і persistence після refresh.}

\taskcard{C13}{Realtime content}{2}{Передавати зміни без reload.}
{Після API-операції hub транслює create/update/delete/order. Власна подія не
дублюється, remote update підсвічується. Blocks API і browser flow протестовані.}

\taskcard{C14}{Presence і focus}{2}{Показати блок іншого редактора.}
{focus/blur передають document id, block id та identity. Group isolation не
дозволяє подіям перейти в інший документ.}

\taskcard{C15}{File API}{3}{Застосувати File, DataTransfer, FileReader.}
{Drop читає DataTransfer.files, другий DataTransfer наповнює native input,
FileReader дає preview. Сервер незалежно перевіряє size/extension. Реальний
text upload пройшов.}

\taskcard{C16}{OAuth 2.0}{3}{Додати зовнішній sign-in.}
{Google handler реєструється лише з ClientId/Secret, callback:
\code{/signin-google}. Для автономного показу є development cookie switcher.
Реальний consent потребує Google credentials і HTTPS callback.}

\taskcard{C18}{Realtime shared board}{4}{Створити спільне малювання.}
{BoardStateService тримає bounded history, SignalR передає draw/clear, новий
клієнт отримує snapshot. Локально перевірено підключення.}

\taskcard{C19}{Cloud/local storage}{3}{Абстрагувати attachments.}
{AttachmentStorage обирає Azure Blob за connection string або local uploads.
БД зберігає URL та metadata. Local path перевірено фактичним upload; Azure
потребує storage account.}

\taskcard{C21}{SSE-помічник}{5}{Потоково передавати прогрес і відповідь.}
{Assistant повертає planning, tool, token, done і використовує tenant context.
Smoke test отримав 200 та done=true; UI показав agent flow.}

\taskcard{C22}{MCP server}{5}{Надати стандартні tools.}
{Окремий stdio server підтримує initialize, discovery, tools/list і tools/call.
Smoke tools/list повернув list\_documents, search\_documents, get\_document.}

\subsection{BackEnd}
\taskcard{B1}{Пагінація}{2}{Обмежити collection response.}
{GET приймає skip/limit, повертає items, total і nextLink; limit 1--100.
Smoke test: total=1, items=1; validation і nextLink покриті тестами.}

\taskcard{B2}{Webhook/Telegram}{2}{Надсилати publish notification.}
{WebhookNotificationService підтримує generic webhook і Telegram Bot API,
не ламаючи CRUD при outage. Код збирається; доставка потребує token/chat id
або test webhook.}

\taskcard{B4}{Full-text search}{2}{Реалізувати релевантний пошук.}
{DocumentSearchService будує tenant-scoped Lucene index і scored results для
AJAX autocomplete. Smoke знайшов 1 результат; UI покрито E2E.}

\taskcard{B5}{Ролі}{3}{Розділити Owner, Editor і Viewer.}
{TenantAccessService перевіряє membership і minimum role. Viewer читає,
Editor змінює контент, Owner керує простором, учасниками та планом.}

\taskcard{B6}{Multi-tenancy і plans}{3}{Ізолювати команди й feature gates.}
{UI/API/search/map/cache/SSE обмежені membership. Free: 2 MB; Premium: 10 MB,
board, assistant, advanced insights. Unit/auth tests доводять правила.}

\taskcard{B7}{Azure Functions}{3}{Створити serverless API.}
{DigestFunctions має anonymous health і function-key digest. Проєкт збирається.
func CLI відсутній, тому local host/cloud runtime потребує окремої перевірки.}

\taskcard{B8}{RabbitMQ worker}{4}{Винести фонову обробку.}
{Producer публікує durable messages, Worker має prefetch=1, ACK/reject і
reconnect. Обидва проєкти збираються; broker runtime виконується окремо через
Docker Compose.}

\taskcard{B9}{API cache}{3}{Кешувати повторне читання.}
{GET by id використовує IMemoryCache на 2 хвилини. Перший response має MISS,
другий HIT; PUT/DELETE інвалідують entry. Smoke і integration test пройшли.}

\subsection{FrontEnd}
\taskcard{F1}{Chart.js}{3}{Візуалізувати статистику.}
{Insights агрегує актуальні DB values й передає dataset у doughnut chart.
Локально видно 1 workspace, 1 document, 3 blocks і published status.}

\taskcard{F2}{jQuery autocomplete}{2}{Шукати без reload.}
{Після 3 символів і 250 ms debounce викликається search API, попередній
request скасовується, результати є links. Перевірено Playwright.}

\taskcard{F3}{Drag-and-drop}{2}{Змінювати порядок блоків.}
{Native drag events формують масив id; server транзакційно перевіряє блоки,
оновлює Position і надсилає порядок SignalR.}

\taskcard{F5}{React SPA}{3}{Створити окремий TypeScript-клієнт.}
{React/Vite SPA має list, search, edit, loading/error/empty states. pnpm build
успішно створив bundle: JS приблизно 223 kB (69.8 kB gzip), CSS 3.2 kB.}

\taskcard{F6}{Pointer Events і Canvas}{5}{Підтримати mouse, pen і touch.}
{Canvas слухає pointerdown/move/up, зберігає normalized coordinates, colour і
width; resize не зміщує strokes, completed stroke передається hub.}

\subsection{QA}
\taskcard{QA1}{API integration tests}{2}{Автоматизувати REST lifecycle.}
{DocumentsApiTests перевіряє list/create/read/update/delete, paging, validation,
not found, conflict і cache. Усі integration tests пройшли.}

\taskcard{QA2}{Browser E2E}{3}{Перевірити справжні user flows.}
{Playwright покриває landing, Lucene autocomplete, block editor, map marker,
robots.txt і sitemap.xml. Фінальний результат: 5/5 passed.}

\taskcard{QA3}{Authorization tests}{3}{Довести tenant isolation.}
{Дві test identities показують: owner бачить seed workspace, non-member не
може list/open. FeatureGateTests покриває Free/Premium.}

\section{REST API і ручна перевірка}
\begin{longtable}{p{0.13\textwidth} p{0.39\textwidth} p{0.35\textwidth}}
\toprule Method & URL & Результат \\
\midrule\endfirsthead
\toprule Method & URL & Результат \\
\midrule\endhead
GET & /api/documents & Paginated summaries \\
GET & /api/documents/\{id\} & Detail + X-Cache \\
POST & /api/documents & 201 Created + Location \\
PUT & /api/documents/\{id\} & Update + invalidation \\
DELETE & /api/documents/\{id\} & 204 No Content \\
GET & /api/search/documents & Lucene results \\
GET & /api/map/workspaces & Markers \\
GET & /api/connectivity/state & Pulse \\
GET & /api/connectivity/long-poll & Wait for version \\
GET & /api/connectivity/websocket & WebSocket upgrade \\
GET & /api/assistant/stream & SSE events \\
\bottomrule\end{longtable}

Ручний smoke test на живому сервері: collection total=1/items=1; два GET одного
документа дали MISS, потім HIT; POST повернув 201; search повернув 1 result;
DELETE повернув 204; Map повернув 1 marker; SSE повернув 200 і done=true.

\begin{lstlisting}[language={},caption={Paginated response}]
{"items":[...],"skip":0,"limit":20,"total":1,"nextLink":null}
\end{lstlisting}

\section{Безпека і конфігурація}
Cookie authentication є базовою. Development switcher дає дві identity без
секретів. Google OAuth активується лише з credentials. Identity зіставляється
з WorkspaceMember; роль не приймається з browser request без перевірки.

\begin{itemize}
\item server-side validation доповнює client validation;
\item anti-forgery захищає MVC state-changing forms;
\item DTO запобігають mass assignment;
\item tenant filters діють для UI, API, search, map, cache і SSE;
\item upload має allowlist і plan-dependent limit;
\item секрети не записані до Git;
\item MCP stdout містить лише protocol messages.
\end{itemize}

\begin{table}[H]\centering\small
\begin{tabularx}{\textwidth}{p{0.47\textwidth} X}
\toprule Параметр & Призначення \\
\midrule
Authentication: Google ClientId / ClientSecret & Google OAuth \\
Application Insights connection string & Azure telemetry \\
Storage: BlobConnectionString & Azure Blob \\
Notifications: Telegram token / ChatId & Telegram \\
Notifications: WebhookUrl & Generic webhook \\
RabbitMq: Host & Broker \\
Assistant: Endpoint / ApiKey / Model & External AI \\
\bottomrule\end{tabularx}\caption{Секретна зовнішня конфігурація}\end{table}

\section{Тестування та результати перевірки}
Перевірку виконано 26 жовтня 2026 року у Windows, .NET SDK 10.0.402,
Node.js 24.22.1, pnpm 11.25.0. Сервер працював на
\url{http://localhost:5036}; browser tests виконувалися у Chromium Playwright.

Рівні: \textbf{6 unit}, \textbf{14 integration}, \textbf{5 E2E}; разом
\textbf{25 автоматизованих тестів}, усі успішні.

\begin{longtable}{>{\bfseries}p{0.30\textwidth} p{0.18\textwidth} p{0.39\textwidth}}
\toprule Перевірка & Результат & Примітка \\
\midrule\endfirsthead
\toprule Перевірка & Результат & Примітка \\
\midrule\endhead
dotnet format & PASS & Без змін \\
dotnet build & PASS & 0 warnings, 0 errors \\
Unit tests & 6/6 PASS & Домен і gates \\
Integration tests & 14/14 PASS & API та authorization \\
React production build & PASS & TypeScript + Vite \\
Playwright E2E & 5/5 PASS & Chromium \\
API CRUD/paging/cache/search & PASS & 201, 204, MISS/HIT \\
Map/SSE/MCP & PASS & marker, done, 3 tools \\
File upload & PASS & Файл потім видалено \\
\bottomrule\end{longtable}

\begin{lstlisting}[language=bash,caption={Основні команди перевірки}]
dotnet restore StudySpace.slnx
dotnet format StudySpace.slnx --verify-no-changes --no-restore
dotnet build StudySpace.slnx --no-restore
dotnet test tests/StudySpace.UnitTests --no-build
dotnet test tests/StudySpace.IntegrationTests --no-build
cd src/studyspace-spa && pnpm build
cd ../../tests/e2e && pnpm test
\end{lstlisting}

\subsection{MCP smoke test}
\begin{lstlisting}[language={},caption={PowerShell tools/list}]
'{"jsonrpc":"2.0","id":1,"method":"tools/list","params":{}}' |
  dotnet run --project src/StudySpace.Mcp --no-build
\end{lstlisting}
Відповідь містила list\_documents, search\_documents і get\_document.

\section{Запуск кожного застосунку}
\subsection{Web}
\begin{lstlisting}[language=bash]
dotnet restore StudySpace.slnx
dotnet run --project src/StudySpace.Web --urls http://localhost:5036
\end{lstlisting}

\subsection{React SPA}
\begin{lstlisting}[language=bash]
cd src/studyspace-spa
pnpm install
pnpm dev
\end{lstlisting}
API має працювати паралельно.

\subsection{Container stack}
\begin{lstlisting}[language=bash]
docker compose up --build
\end{lstlisting}
Після інсталяції Docker: web використовує порт 8080, RabbitMQ використовує
порт 5672, management UI використовує порт 15672.

\subsection{Azure Functions}
\begin{lstlisting}[language=bash]
cd src/StudySpace.Functions
func start
\end{lstlisting}
Потрібні Functions Core Tools і StudySpaceApiUrl.

\subsection{MCP}
\begin{lstlisting}[language=bash]
dotnet run --project src/StudySpace.Mcp
\end{lstlisting}
Це не сайт: stdin/stdout підключаються до MCP-клієнта.

\section{Автоматизація перевірок}
GitHub Actions автоматично виконує форматування, збірку, .NET-тести,
production-збірку frontend, Playwright E2E та Docker image build у чистому
Ubuntu environment.

\section{Висновки}
Реалізовано StudySpace і 36 пунктів на 104 бали: один монорепозиторій,
5 запускних застосунків і 10 технічних проєктів. Пов'язані MVC, API, data,
auth, search, files і realtime-функції об'єднано, а React, Functions, Worker
і MCP залишено окремими процесами.

Перевірка підтвердила збірку без warnings/errors, 25 тестів, React bundle, CRUD,
pagination, cache MISS/HIT, Lucene, map, SSE і MCP tools/list. Роботу
зафіксовано 12 актуальними скріншотами. Рішення має завершену структуру та
може бути розгорнуте у зовнішньому середовищі після налаштування відповідних
сервісів.
\end{document}
